\sectionkicker{Source-backed technical notes}
\subsection*{モデルを動かす技術が競争軸になった――QLoRA、FlashAttention-2、vLLM: Technical Notes}
本欄は記事本文で圧縮した一次資料上の情報を、比較・再検証しやすい形で整理したものである。時系列、確認済みの主張、留意点、一次資料URLを掲載する。
\medskip
\subsection*{Theme at a glance}
\addcontentsline{toc}{subsection}{Theme at a glance}
\begin{center}
\footnotesize
\begin{tabularx}{\linewidth}{@{}p{0.20\linewidth}p{0.10\linewidth}p{0.14\linewidth}X@{}}
\toprule
資料 & 位置づけ & 種別 & 時系列 \\
\midrule
vLLM / PagedAttention & 主要資料 & Framework & — \\
QLoRA & 主要資料 & 論文 & — \\
FlashAttention-2 & 主要資料 & 論文 & — \\
Phi-2 & 補足資料 & モデル & — \\
\bottomrule
\end{tabularx}
\end{center}
\normalsize

\begin{technicalnote}{vLLM / PagedAttention}{主要資料}
\begingroup
% reader-facing Technical Notes break policy
\widowpenalty=10000
\clubpenalty=10000
\displaywidowpenalty=10000
\begin{tabularx}{\linewidth}{@{}>{\bfseries}p{0.22\linewidth}X@{}}
組織 & vLLM Project \\
種別 & Framework \\
時系列 & — \\
\end{tabularx}
\smallskip
{\bfseries 一次資料から整理したtechnical points}
\begin{itemize}[leftmargin=1.5em,itemsep=0.35em]
\item \textbf{一次情報で確認できる事実}: 一次資料で「vLLM / PagedAttention」の2023年における公開・リリースの経緯を確認できる。能力や性能に関する評価は、提供元・プロジェクト・著者の主張として扱う。
\end{itemize}
\begin{minipage}{\linewidth}
% reader-facing Technical Notes generic boundary/source tail group
{\bfseries 読む際の境界}
\begin{itemize}[leftmargin=1.5em,itemsep=0.35em]
\item \textbf{分析上の留意点}: 一次資料で確認できる範囲の事実を記録しており、提供元・プロジェクト・著者による評価を独立再現済みの結果として扱わない。
\end{itemize}
\begin{samepage}
{\bfseries 一次資料}
\begingroup\sloppy
\Urlmuskip=0mu plus 2mu\relax
\begin{itemize}[leftmargin=1.5em,itemsep=0.25em]
\item {\scriptsize\url{http://arxiv.org/abs/2309.06180v1}}
\item {\scriptsize\url{https://github.com/vllm-project/vllm/releases/tag/v0.1.3}}
\item {\scriptsize\url{https://vllm.ai/blog/2023-06-20-vllm}}
\end{itemize}
\endgroup
\end{samepage}
\end{minipage}
\endgroup
\end{technicalnote}
\medskip

\begin{technicalnote}{QLoRA}{主要資料}
\begingroup
% reader-facing Technical Notes break policy
\widowpenalty=10000
\clubpenalty=10000
\displaywidowpenalty=10000
\begin{tabularx}{\linewidth}{@{}>{\bfseries}p{0.22\linewidth}X@{}}
組織 & authors \\
種別 & 論文 \\
時系列 & — \\
\end{tabularx}
\smallskip
{\bfseries 一次資料から整理したtechnical points}
\begin{itemize}[leftmargin=1.5em,itemsep=0.35em]
\item \textbf{Author claim}: 一次資料で「QLoRA」が2023年の技術動向として記録されていることを確認できる。技術内容や評価結果は、原著者による主張として扱い、独立再現済みの結果とはみなさない。
\end{itemize}
\begin{minipage}{\linewidth}
% reader-facing Technical Notes generic boundary/source tail group
{\bfseries 読む際の境界}
\begin{itemize}[leftmargin=1.5em,itemsep=0.35em]
\item \textbf{分析上の留意点}: 一次資料で確認できる範囲の事実を記録しており、提供元・プロジェクト・著者による評価を独立再現済みの結果として扱わない。
\end{itemize}
\begin{samepage}
{\bfseries 一次資料}
\begingroup\sloppy
\Urlmuskip=0mu plus 2mu\relax
\begin{itemize}[leftmargin=1.5em,itemsep=0.25em]
\item {\scriptsize\url{http://arxiv.org/abs/2305.14314v1}}
\end{itemize}
\endgroup
\end{samepage}
\end{minipage}
\endgroup
\end{technicalnote}
\medskip

\begin{technicalnote}{FlashAttention-2}{主要資料}
\begingroup
% reader-facing Technical Notes break policy
\widowpenalty=10000
\clubpenalty=10000
\displaywidowpenalty=10000
\begin{tabularx}{\linewidth}{@{}>{\bfseries}p{0.22\linewidth}X@{}}
組織 & authors \\
種別 & 論文 \\
時系列 & — \\
\end{tabularx}
\smallskip
{\bfseries 一次資料から整理したtechnical points}
\begin{itemize}[leftmargin=1.5em,itemsep=0.35em]
\item \textbf{Author claim}: 一次資料で「FlashAttention-2」が2023年の技術動向として記録されていることを確認できる。技術内容や評価結果は、原著者による主張として扱い、独立再現済みの結果とはみなさない。
\end{itemize}
\begin{minipage}{\linewidth}
% reader-facing Technical Notes generic boundary/source tail group
{\bfseries 読む際の境界}
\begin{itemize}[leftmargin=1.5em,itemsep=0.35em]
\item \textbf{分析上の留意点}: 一次資料で確認できる範囲の事実を記録しており、提供元・プロジェクト・著者による評価を独立再現済みの結果として扱わない。
\end{itemize}
\begin{samepage}
{\bfseries 一次資料}
\begingroup\sloppy
\Urlmuskip=0mu plus 2mu\relax
\begin{itemize}[leftmargin=1.5em,itemsep=0.25em]
\item {\scriptsize\url{http://arxiv.org/abs/2307.08691v1}}
\end{itemize}
\endgroup
\end{samepage}
\end{minipage}
\endgroup
\end{technicalnote}
\medskip

\begin{technicalnote}{Phi-2}{補足資料}
\begingroup
% reader-facing Technical Notes break policy
\widowpenalty=10000
\clubpenalty=10000
\displaywidowpenalty=10000
\begin{tabularx}{\linewidth}{@{}>{\bfseries}p{0.22\linewidth}X@{}}
組織 & Microsoft Research \\
種別 & モデル \\
時系列 & — \\
\end{tabularx}
\smallskip
{\bfseries 一次資料から整理したtechnical points}
\begin{itemize}[leftmargin=1.5em,itemsep=0.35em]
\item \textbf{一次情報で確認できる事実}: 一次資料で「Phi-2」の2023年における公開・リリースの経緯を確認できる。能力や性能に関する評価は、提供元・プロジェクト・著者の主張として扱う。
\end{itemize}
\begin{minipage}{\linewidth}
% reader-facing Technical Notes generic boundary/source tail group
{\bfseries 読む際の境界}
\begin{itemize}[leftmargin=1.5em,itemsep=0.35em]
\item \textbf{分析上の留意点}: 一次資料で確認できる範囲の事実を記録しており、提供元・プロジェクト・著者による評価を独立再現済みの結果として扱わない。
\end{itemize}
\begin{samepage}
{\bfseries 一次資料}
\begingroup\sloppy
\Urlmuskip=0mu plus 2mu\relax
\begin{itemize}[leftmargin=1.5em,itemsep=0.25em]
\item {\scriptsize\url{https://www.microsoft.com/en-us/research/blog/phi-2-the-surprising-power-of-small-language-models/}}
\end{itemize}
\endgroup
\end{samepage}
\end{minipage}
\endgroup
\end{technicalnote}
\medskip
